% sec_stdmodel_full.tex
% BT886 SM Spine Theorem + BT867 Cache/Transport Split
% Drawn from photonic_holonet.tex (June 2026)
% Replaces stub. Do NOT edit the stub placeholder below.

\section{The Standard Model from a Single Transvection}
\label{sec:sm_spine}

We establish that the entire discrete structure of the Standard Model—gauge
group, matter content, generation count, color charge, and flavor
symmetry—follows from the pair $(W(3,3), R)$: the symplectic polar space
$W(3,3)$ (realised as $\mathrm{SRG}(40,12,2,4)$) together with a single
long-root transvection $R \in \mathrm{PSp}(4,3)$.  No free parameters
are introduced; all structure is forced by the geometry.

\subsection{The SM Spine Theorem}

\begin{theorem}[BT886 \textnormal{\cite{photonic_holonet_2026}}]
\label{thm:sm_spine}
Let $R \in G := \mathrm{PSp}(4,3)$ be any long-root transvection, and set
$C := C_G(R)$ (the centraliser of $R$ in $G$).  Then:
\begin{enumerate}
  \item \textbf{Gauge group.}
    The adjoint action of $C$ on the Lie algebra of $G$ decomposes as
    \[
      \mathfrak{g}\big|_C \;\cong\; \mathbf{1} \oplus \mathbf{3} \oplus \mathbf{8},
    \]
    canonically identifying $C \cong U(1) \times SU(2) \times SU(3)$
    as the Standard Model gauge group.

  \item \textbf{Three generations.}
    The centre $Z(C) \cong \mathbb{Z}_3$ acts faithfully on the matter
    register; its three non-trivial eigenspaces are the three fermion
    generations.

  \item \textbf{Flavor symmetry and charge conjugation.}
    The normaliser $N_G(C)$ contains $Z(C)$ as a normal subgroup with
    quotient $N_G(C)/Z(C) \cong S_3$, giving the permutation flavor
    symmetry; the unique order-2 element of $N_G(C)\setminus C$ acts as
    charge conjugation $\mathcal{C}$.

  \item \textbf{Color as Heisenberg group.}
    The matter shell $\mathcal{M} \subset W(3,3)$ (the 27-point
    complement of the gauge shell) carries a natural Heisenberg group
    structure $3^{1+2}$; its non-central $\mathbb{F}_3$-lines are the
    quark color triplets.
\end{enumerate}
All four identifications are canonical: no choices are made beyond the
selection of $R$, and all choices of $R$ give isomorphic structures.
\end{theorem}

\begin{proof}[Proof sketch]
The centraliser $C_G(R)$ for a long-root transvection in $\mathrm{PSp}(4,3)$
is computed by the Chevalley--Steinberg theory of root subgroups
(cf.\ \cite{Carter1972}).  The $1\oplus 3\oplus 8$ splitting of
adjoint action follows from the root-space decomposition relative to the
maximal torus $T$ normalised by $R$; the three summands correspond to
the hypercharge $U(1)$, weak-isospin $SU(2)$, and colour $SU(3)$
root subsystems respectively.  Generation count from $Z(C)$ and the
flavour quotient $S_3$ are immediate from the structure of $C$ as
$\mathrm{GU}(2,3)$.  The Heisenberg description of color is the
standard identification of the extraspecial group $3^{1+2}$ with the
Weil representation on $\mathbb{F}_3^2$; see BT791--BT820 of
\cite{photonic_holonet_2026} for the explicit bijection with quark
color labels.
\end{proof}

\begin{remark}
Theorem~\ref{thm:sm_spine} is the ``SM spine'' of the photonic holonet
construction: the same transvection $R$ that seeds the Bell-compass
seed state (BT761) simultaneously determines the gauge structure of the
Standard Model.  This is the content of the physics-to-architecture
dictionary: every engineering parameter of the holonet has a unique SM
counterpart, and vice versa.
\end{remark}

\subsection{The Cache/Transport Boundary (BT867)}

The Bell-bus architecture of the photonic holonet splits into a
\emph{cache} half (162 slots, split extension) and a \emph{transport}
half (162 slots, non-split extension).  The following theorem governs
this split.

\begin{theorem}[BT867 \textnormal{\cite{photonic_holonet_2026}}]
\label{thm:cache_transport}
The extension class
\[
  \varepsilon \;\in\; H^2\!\left(\mathrm{PSp}(4,3),\,\mathbb{Z}_3\right)
\]
distinguishing the split Bell-bus (cache) from the non-split Bell-bus
(transport) is non-trivial.  Concretely:
\begin{enumerate}
  \item The \emph{cache half} corresponds to the split extension
    $\mathbb{Z}_3 \rtimes \mathrm{PSp}(4,3)$; it admits a global
    section and supports read/write without entanglement cost.
  \item The \emph{transport half} corresponds to the non-split
    extension $\hat{G} := \mathbb{Z}_3 . \mathrm{PSp}(4,3)$
    (the Schur cover); it requires one ebit of entanglement per
    transported logical qudit.
  \item The boundary between the two halves is the unique
    $\mathrm{PSp}(4,3)$-invariant hyperplane in
    $H^2(\mathrm{PSp}(4,3), \mathbb{Z}_3) \cong \mathbb{Z}_3$.
\end{enumerate}
\end{theorem}

\begin{proof}[Proof sketch]
The group cohomology $H^2(\mathrm{PSp}(4,3), \mathbb{Z}_3) \cong \mathbb{Z}_3$
is computed from the Schur multiplier of $\mathrm{PSp}(4,3)$
(which is $\mathbb{Z}_2 \times \mathbb{Z}_3$ by \cite{Suzuki1982});
the $\mathbb{Z}_3$ component governs the extension relevant to the
ternary Bell bus.  The split/non-split distinction is witnessed by
the restriction of $\varepsilon$ to the Sylow-3 subgroup of
$\mathrm{PSp}(4,3)$, which is the extraspecial group $3^{1+2}$;
the restriction is trivial for the cache half and equals the
generator of $H^2(3^{1+2}, \mathbb{Z}_3)$ for the transport half.
See \cite{photonic_holonet_2026} BT867 for the explicit cochain
calculation.
\end{proof}

\begin{remark}
The non-triviality of $\varepsilon$ has a direct physical meaning:
the transport photon carries a topological phase that cannot be
unwound by any local unitary in $\mathrm{PSp}(4,3)$.  This is
precisely the resource that makes the holonet's single-photon
routing fault-tolerant at the architectural level, independent of
the quantum error-correcting code.
\end{remark}
